%%%%%%%%%%%%Outline%%%%%%%%%%%%%%%%%%

%%% Section 8 Related Work

%%%
% message<bytecode optimizers stay within the instruction set, \tool grows the operation set>
% bytecode optimizers improve eBPF programs within the existing instruction set
% K2 searches for shorter bytecode with stochastic synthesis
% Merlin rewrites at the LLVM-IR and bytecode levels
% EPSO applies rewrite rules from offline enumerative superoptimization
% ePass lifts bytecode into SSA form
% KFuse merges chains of eBPF programs after verification
% \tool grows the set of operations available to a program

%%%
% message<the existing routes to new operations hide semantics or revise the instruction set>
% the existing routes to new operations are kfuncs and instruction-set extensions
% a kfunc exposes a kernel function for eBPF programs to call
% the verifier checks the arguments, the body runs outside the verifier's analysis
% hXDP extends the instruction set for FPGAs, translated programs run outside the verifier's coverage
% the eBPF instruction set is an IETF standard, a revision is a standardization effort
% a \tool operation shows the verifier its full semantics via the proof sequence
% a new operation ships as a module plus a recognizer pattern, no instruction-set revision

%%%
% message<other systems move execution or re-found safety, \tool keeps both in place>
% another line moves where eBPF runs or what its safety rests on
% bpftime runs eBPF in userspace
% eBPF for Windows checks with PREVAIL in userspace, loads signed driver modules
% Craun et al. verify, compile, and sign on a separate host for embedded devices
% KFlex moves only the extension's own memory safety and termination to runtime checks
% Rex drops the verifier and accepts extensions written in safe Rust
% MOAT and SafeBPF add runtime isolation beneath the static checks
% with \tool, programs stay in the kernel under the existing verifier's load-time check
% the one addition to the trusted computing base is the native emit

%%%
% message<verified compilation targets fixed translators, \lang proves per operation>
% verified compilation provides the techniques behind the \lang proofs
% CompCert established mechanized semantic preservation for a realistic compiler
% Alive2 checks refinement for each LLVM transformation
% Jitk compiles classic BPF through a Coq-verified JIT replacing the seccomp interpreter
% Jitterbug proves per-instruction eBPF JIT equivalence, JitSynth synthesizes such JITs
% CertrBPF derives a verified C interpreter for eBPF from a Coq specification
% all of these systems target a fixed instruction set
% each \lang operation carries a Lean 4 proof that emit and proof sequence agree

%%%
% message<where machine-checked proofs sit relative to the kernel>
% beyond the translators, proofs also target eBPF's verification and loading steps
% Agni verifies offline the range analysis on whose soundness \tool's safety argument rests
% BCF supplies proofs from userspace, the kernel checks the proofs at load time
% \lang's equivalence proofs stay offline as development-time evidence
% proof-carrying code is the closest ancestor of \tool
% both admit untrusted code into the kernel only after an in-kernel check
% PCC installs a new proof checker in the kernel
% \tool reuses the existing eBPF verifier as the checker

%%%%%%%%%%%%Outline%%%%%%%%%%%%%%%%%%

\section{Related Work}
\label{sec:related}
% 相关工作

\noindent\textbf{eBPF bytecode optimization.}
% \noindent\textbf{eBPF 字节码优化。}
Bytecode optimizers improve eBPF programs within the existing instruction set.
% 字节码优化器在现有指令集内改进 eBPF 程序。
K2~\cite{xu2021synthesizing} searches for shorter bytecode with stochastic synthesis.
% K2 使用随机合成搜索更短的字节码。
Merlin~\cite{mao2024merlin} rewrites programs at the LLVM-IR and bytecode levels.
% Merlin 在 LLVM-IR 和字节码级别重写程序。
EPSO~\cite{zhu2025epso} applies rewrite rules from offline enumerative superoptimization.
% EPSO 应用来自离线枚举超优化的重写规则。
ePass~\cite{xiang2025epass} lifts bytecode into SSA form for program transformation.
% ePass 将字节码提升为 SSA 形式以进行程序转换。
KFuse~\cite{kfuse} merges chains of eBPF programs into one program after verification.
% KFuse 在验证后将 eBPF 程序链合并为一个程序。
\tool extends the set of operations available to a program.
% \tool 扩展程序可用的操作集。
%
\note{add fun optimization paper}

\noindent\textbf{Extending eBPF's operations.}
% \noindent\textbf{扩展 eBPF 的操作。}
The existing routes to new operations are kfuncs and instruction-set extensions.
% 通往新操作的现有路线是 kfunc 和指令集扩展。
A kfunc exposes a kernel function for eBPF programs to call~\cite{kfuncs}.
% kfunc 暴露一个内核函数供 eBPF 程序调用。
The verifier checks a kfunc call's arguments against BTF types, and the body runs as native code outside the verifier's analysis.
% 验证器根据 BTF 类型检查 kfunc 调用的参数，函数体作为本地代码在验证器分析之外运行。
hXDP~\cite{brunella2020hxdp} extends the eBPF instruction set for FPGA targets, and the translated programs run outside the verifier's coverage.
% hXDP 为 FPGA 目标扩展 eBPF 指令集，翻译后的程序在验证器覆盖范围之外运行。
The eBPF instruction set is now an IETF standard~\cite{rfc9669}, so a revision is a standardization effort.
% eBPF 指令集现在是 IETF 标准，因此修订是一项标准化工作。
A \tool operation shows the verifier its full semantics via a proof sequence in the lowered program (\S\ref{subsec:pipeline}).
% \tool 操作通过降低程序中的证明序列向验证器展示其完整语义。
A new operation ships as a kernel module plus a recognizer pattern, with no instruction-set revision (\S\ref{subsec:adding}).
% 新操作作为内核模块加识别器模式发布，无需指令集修订。

\noindent\textbf{Relocating execution or the safety argument.}
% \noindent\textbf{重定位执行或安全论证。}
Another line of work relocates eBPF execution or changes the foundation of its safety argument.
% 另一条工作线重定位 eBPF 执行或改变其安全论证的基础。
bpftime~\cite{zheng2025extending} runs eBPF programs in userspace.
% bpftime 在用户空间运行 eBPF 程序。
eBPF for Windows~\cite{windows-ebpf} checks programs in userspace with the PREVAIL verifier~\cite{gershuni2019simple} and loads checked programs as signed driver modules.
% eBPF for Windows 使用 PREVAIL 验证器在用户空间检查程序，并将检查过的程序作为签名驱动模块加载。
Craun et al.~\cite{craun2023enabling} verify, compile, and sign programs on a separate host for embedded devices.
% Craun 等人在单独主机上为嵌入式设备验证、编译和签名程序。
KFlex~\cite{dwivedi2024fast} moves only the extension's own memory safety and termination to runtime checks.
% KFlex 仅将扩展自身的内存安全和终止性移至运行时检查。
Rex~\cite{jia2025rex} drops the verifier and accepts extensions written in safe Rust.
% Rex 放弃验证器，接受用安全 Rust 编写的扩展。
MOAT~\cite{lu2024moat} and SafeBPF~\cite{soo2024safebpf} add runtime isolation beneath the verifier's static checks.
% MOAT 和 SafeBPF 在验证器的静态检查之下添加运行时隔离。
With \tool, programs stay in the kernel, under the existing verifier's load-time check.
% 使用 \tool，程序留在内核中，在现有验证器的加载时检查下。
Beyond \tool's kernel patch, the only per-operation addition to the TCB is the native emit (\S\ref{subsec:safety}).
% 除了 \tool 的内核补丁，每个操作对 TCB 的唯一增加是本地发射。

\noindent\textbf{Verified compilation and proof-carrying code.}
% \noindent\textbf{经验证的编译和携带证明的代码。}
Verified compilation provides the techniques behind the \lang proofs.
% 经验证的编译提供 \lang 证明背后的技术。
CompCert~\cite{leroy2009formal} established mechanized semantic preservation for a realistic compiler.
% CompCert 为现实编译器建立了机械化的语义保持。
Alive2~\cite{lopes2021alive2} checks refinement for each LLVM transformation.
% Alive2 为每个 LLVM 转换检查精化。
Jitk~\cite{wang2014jitk} compiles classic BPF through a Coq-verified JIT that replaces the seccomp interpreter path.
% Jitk 通过 Coq 验证的 JIT 编译经典 BPF，替换 seccomp 解释器路径。
Jitterbug~\cite{nelson2020specification} proves behavioral equivalence for per-instruction eBPF JIT translation with an SMT solver, and JitSynth~\cite{vangeffen2020jitsynth} synthesizes such verified JITs from interpreter semantics.
% Jitterbug 使用 SMT 求解器证明逐指令 eBPF JIT 翻译的行为等价，JitSynth 从解释器语义合成这样的验证 JIT。
CertrBPF~\cite{yuan2022certrbpf} derives a verified C implementation of an eBPF interpreter from a Coq specification.
% CertrBPF 从 Coq 规范推导出 eBPF 解释器的经验证 C 实现。
All of these systems target a fixed instruction set.
% 所有这些系统都针对固定的指令集。
Each \lang operation instead includes its own Lean 4 proof that the native emit computes the same result as its proof sequence (\S\ref{sec:formal-verification}).
% 每个 \lang 操作则包含其自己的 Lean 4 证明，证明本地发射与其证明序列计算相同的结果。

Beyond the translators, machine-checked proofs also target eBPF's verification and loading steps.
% 除了翻译器，机器检查的证明还针对 eBPF 的验证和加载步骤。
Agni~\cite{vishwanathan2023verifying} verifies offline the eBPF verifier's range analysis, on whose soundness \tool's safety argument rests.
% Agni 离线验证 eBPF 验证器的范围分析，\tool 的安全论证依赖于其正确性。
BCF~\cite{sun2025prove} supplies abstraction-refinement proofs from userspace, which the kernel checks at load time.
% BCF 从用户空间提供抽象精化证明，内核在加载时检查。
\lang's equivalence proofs stay offline, as development-time evidence that never enters the kernel.
% \lang 的等价证明保持离线，作为永不进入内核的开发时证据。
Proof-carrying code (PCC)~\cite{necula1997pcc} is the closest ancestor of \tool.
% 携带证明的代码（PCC）是 \tool 最接近的祖先。
Both admit untrusted code into the kernel only after an in-kernel check, but whereas PCC installs a new proof checker, \tool reuses the existing eBPF verifier.
% 两者都只在内核内检查后才将不可信代码纳入内核，但 PCC 安装新的证明检查器，而 \tool 重用现有的 eBPF 验证器。